\DocumentMetadata{lang=es-DO,testphase={phase-III,math}}
\documentclass[11pt]{scrartcl}
\usepackage{icv}

% -----------------------------------------------------------------------
% Migrado de legacy/Acueductos y Alcantarillados/Proyecto/05_rubrica.tex
% -- rúbrica real completa. La sección "Proyecto escrito" ya usaba en el
% original los mismos 4 niveles fijos que icvrubric adopta (Excelente/
% Satisfactorio/En desarrollo/Insuficiente) con peso% por criterio, así
% que migra 1:1. "Presentación oral" tiene una forma distinta (columna
% única "Se espera", no 4 niveles) -- se deja como tabla propia, no
% forzada a icvrubric.
% -----------------------------------------------------------------------
\icvsetup{
  asignatura  = {ICV 442-T -- Acueductos y Alcantarillados},
  titulocorto = {Rúbrica de evaluación},
  titulo      = {Rúbrica de evaluación},
  subtitulo   = {Proyecto escrito y presentación oral},
  profesor    = {Dr. Ricardo Hernández Moreira},
  periodo     = \icvciclo{2026}{3},
  version     = {v1},
}

\begin{document}
\icvmaketitle

\begin{icvnote}[Escala de evaluación]
\begin{itemize}
  \item \textbf{Excelente}: cumple todos los requisitos con rigor y
    precisión; el resultado es correcto y completo.
  \item \textbf{Satisfactorio}: cumple los requisitos esenciales; puede
    haber errores menores que no afectan el resultado global.
  \item \textbf{En desarrollo}: cumple parcialmente; hay evidencia de
    comprensión del objetivo, pero los errores afectan la calidad del
    resultado.
  \item \textbf{Insuficiente}: no cumple los requisitos mínimos, o el
    entregable está ausente.
\end{itemize}
\end{icvnote}

\section{Proyecto escrito}

\begin{icvrubric}{100 puntos}

\criterio{Calidad de la información básica}{15}
\nivel{Excelente}{Fuentes oficiales claras, proyección bien
  justificada y cálculos trazables.}
\nivel{Satisfactorio}{Información suficiente con algunas debilidades
  de trazabilidad o justificación.}
\nivel{En desarrollo}{Información parcial o poco sustentada.}
\nivel{Insuficiente}{Fuentes débiles o ausencia de base suficiente.}

\criterio{Fundamentos hidráulicos}{35}
\nivel{Excelente}{Supuestos, cálculos y modelos coherentes; resultados
  verificados.}
\nivel{Satisfactorio}{Trabajo correcto con omisiones menores.}
\nivel{En desarrollo}{Inconsistencias parciales o verificación
  incompleta.}
\nivel{Insuficiente}{Errores graves de criterio o falta de
  verificación.}

\criterio{Cumplimiento normativo y justificación técnica}{25}
\nivel{Excelente}{Parámetros clave bien sustentados y decisiones
  defendibles.}
\nivel{Satisfactorio}{Mayormente sustentado, con algunos vacíos.}
\nivel{En desarrollo}{Justificación superficial o dispersa.}
\nivel{Insuficiente}{Criterios adoptados sin sustento suficiente.}

\criterio{Claridad y presentación}{15}
\nivel{Excelente}{Documento ordenado, legible y bien estructurado.}
\nivel{Satisfactorio}{Presentación adecuada con detalles por mejorar.}
\nivel{En desarrollo}{Estructura confusa en varias partes.}
\nivel{Insuficiente}{Documento desordenado o difícil de seguir.}

\criterio{Viabilidad y especificaciones}{10}
\nivel{Excelente}{Materiales, detalles y mantenimiento consistentes
  con el contexto.}
\nivel{Satisfactorio}{En general viable, aunque incompleto en algunos
  puntos.}
\nivel{En desarrollo}{Viabilidad poco discutida o mal cerrada.}
\nivel{Insuficiente}{Sin consideración seria de viabilidad o
  mantenimiento.}

\end{icvrubric}

\section{Presentación oral}

\begin{center}
\begin{tabular}{p{0.34\linewidth} p{0.49\linewidth} p{0.08\linewidth}}
  \toprule
  \textbf{Criterio} & \textbf{Se espera} & \textbf{Peso} \\
  \midrule
  Síntesis y claridad & Exposición ordenada, dentro del tiempo, con
    lenguaje técnico claro. & 30\% \\
  Dominio técnico & Capacidad para explicar criterios, defender
    supuestos y responder preguntas. & 40\% \\
  Material visual & Gráficos, planos y resultados legibles, sin
    sobrecarga. & 15\% \\
  Trabajo en equipo & Participación equilibrada y transiciones
    razonables. & 15\% \\
  \bottomrule
\end{tabular}
\end{center}

\section{Observaciones de uso}
\begin{itemize}
  \item La rúbrica se aplica tanto al producto final como al nivel de
    control técnico demostrado por el equipo.
  \item Un documento muy extenso no compensa la falta de criterio en
    decisiones clave.
  \item La calidad de la presentación oral no sustituye debilidades del
    proyecto escrito, ni a la inversa.
\end{itemize}

\end{document}
